% \section{Limitations}
% \label{saber:sec:limitations}


\section{Conclusion}
\label{saber:sec:conclusion}

In this chapter, we present Saber, a scalable zero-shot framework for reference-to-video generation that eliminates the need for explicitly R2V datasets.
Trained solely on large-scale video-text pairs, Saber leverages a masked training strategy, a tailored attention mechanism, and mask augmentation to achieve identity-consistent, natural, and coherent video generation.
It further scales to multiple references, supporting both multi-identity and multi-view inputs without additional data preparation or changes to the training pipeline.
Extensive experiments on the OpenS2V-Eval benchmark demonstrate that Saber consistently outperforms methods trained on explicit R2V data.
These results show that effective R2V models can be trained without dedicated datasets, paving the way for future research in scalable and generalizable reference-to-video generation.

\paragraph{Limitations.}
While Saber achieves strong zero-shot performance and scalability, several limitations remain.
First, R2V generation may collapse when the number of reference images increases significantly (\eg, $12$), resulting in fragmented compositions where references are combined without coherent understanding.
Second, Saber primarily focuses on identity preservation and visual coherence, while fine-grained motion control and temporal consistency under complex prompts remain challenging.
Future work can explore more effective integration of numerous reference images into unified video generation, as well as adaptive guidance to further improve controllability and realism in reference-to-video generation.
